\section{Exceptional intersections and tangent foliations}
\label{sec:intrinsic}

The quantities $J_{v,e}$ in Lemma~\ref{mixed:flux} measure more than
the number of blowups incident to a factor. They are first moments of
exceptional facets, and they can be recovered from ordinary intersection
numbers on the blowup centres. We establish this interpretation without
assuming that the remaining factors are indexed by a tree. It identifies
the tangent subsheaf whose degree is being estimated and gives a
criterion in terms of a specified collection of exceptional contractions.

\subsection{The radial foliation}

Let $B$ be a smooth projective toric variety and let $b\ge2$. Choose a
torus-fixed point $p\in\PP^b$, distinct invariant hyperplanes
$\mathcal H_1,\ldots,\mathcal H_m$ through $p$, where $0\le m\le b$,
and distinct invariant prime divisors $D_1,\ldots,D_m$ on $B$. Let
\[
 \pi:X\longrightarrow\PP^b\times B
\]
be the successive blowup of the strict transforms of
$\mathcal H_i\times D_i$. Assume that $X$ is Fano, and put
\begin{equation}\label{intrinsic:classes}
 n=\dim X,\qquad H=-K_X,\qquad
 f=\operatorname{pr}_1\circ\pi,\qquad
 L=f^*\mathcal O_{\PP^b}(1).
\end{equation}
The pairs of rays defining the centres are disjoint. In every original
smooth cone their subdivisions act on disjoint pairs of basis vectors.
They therefore commute, all the centres remain smooth, and each
exceptional divisor $E_i$ can be contracted while retaining the other
blowups. The centres themselves need not be disjoint.

Projection from $p$ induces a rational map
\[
 X\dashrightarrow\PP^{b-1}\times B.
\]
Let $\mathcal F$ be the saturation in $T_X$ of the tangent directions
to its general fibres. Thus $\mathcal F$ is the birational transform
of the foliation by lines through $p$, with the $B$-coordinate fixed.

\begin{proposition}\label{intrinsic:foliation}
The sheaf $\mathcal F$ is a rank-one toric foliation and
$\mathcal F\simeq\mathcal O_X(L)$. Its normalized slope is
\begin{equation}\label{intrinsic:slope}
 A:=\frac{\mu_H(\mathcal F)}{\mu_H(T_X)}
   =\frac{nLH^{n-1}}{H^n}
   =\left.\frac{d}{dt}\log\operatorname{vol}(H+tL)\right|_{t=0}.
\end{equation}
In particular, $A\ge1$ excludes stability of $T_X$, and $A>1$ gives
a strict rank-one destabilizer.
\end{proposition}

\begin{proof}
Choose basis rays $r_1,\ldots,r_b$ for the hyperplanes through $p$
and put $r_0=-\sum_jr_j$. On the dense torus, the radial foliation
has tangent line $\CC r_0$. Since $b\ge2$, no other ray of the
projective-space factor lies on this line. Rays of $B$ lie in the
other lattice summand, and the new exceptional rays have nonzero
components in both summands. Thus $r_0$ is the only ray on the line.
The induced tangent filtrations from Lemma~\ref{lem:klyachko} give
$c_1(\mathcal F)=D_{r_0}$. The invariant hyperplane opposite $p$
contains none of the blowup centres, so its strict transform has class
$D_{r_0}=L$.

A saturated rank-one subsheaf of a locally free sheaf on a smooth
variety is reflexive: the quotient is torsion-free, and the depth
lemma gives depth at least two in codimension at least two. It is
therefore a line bundle. This proves the sheaf identification; it
does not assert that the foliation is nonsingular everywhere.
The first equality in \eqref{intrinsic:slope} now follows from
$\mu_H(T_X)=H^n/n$. Near the ample class $H$, algebraic volume is
the top intersection number. Differentiating $(H+tL)^n$ proves the
last equality and the slope conclusions.
\end{proof}

The restriction $b\ge2$ is essential. For a $\PP^1$ factor the two
original rays lie on the same line, and both contribute to its tangent
subsheaf. This is why the leaf estimates in Lemma~\ref{mixed:flux}
are kept separate from the present foliation statement.

\subsection{Mixed intersections on the exceptional centres}

Contract only $E_i$, obtaining $\pi_i:X\to Y_i$ with smooth centre
$Z_i$ of dimension $d=n-2$. The centre is the transverse intersection
of the transforms of $\mathcal H_i\times B$ and $\PP^b\times D_i$.
Their normal line bundles on $Z_i$, in this order, will be denoted
$N_{i,A}$ and $N_{i,B}$. Thus
\[
 N_{Z_i/Y_i}=N_{i,A}\oplus N_{i,B},\qquad
 q_i:E_i=\PP_{\rm lines}(N_{i,A}\oplus N_{i,B})\longrightarrow Z_i.
\]
Let $S_{i,0}$ be the section corresponding to the normal direction
$N_{i,B}$, and let $S_{i,1}$ correspond to $N_{i,A}$. The first is
the intersection with the strict transform of $\mathcal H_i\times B$.
Identify the sections with $Z_i$ and set
\begin{equation}\label{intrinsic:polarizations}
 M_{i,0}=H|_{S_{i,0}}=-K_{Z_i}+N_{i,A},\qquad
 M_{i,1}=H|_{S_{i,1}}=-K_{Z_i}+N_{i,B}.
\end{equation}
These are ample classes. The displayed equalities use additive notation
for line bundles. Indeed, if $\xi=c_1(\mathcal O_{\PP_{\rm lines}(N)}(1))$,
then
\[
 E_i|_{E_i}=-\xi,\qquad
 H|_{E_i}=q_i^*(-K_{Y_i}|_{Z_i})+\xi,
\]
while $\xi$ restricts to $-N_{i,B}$ on $S_{i,0}$ and to $-N_{i,A}$
on $S_{i,1}$. Adjunction gives \eqref{intrinsic:polarizations}.

For the next proposition all products of the $M_{i,j}$ are top
intersection numbers on $Z_i$. Define
\begin{equation}\label{intrinsic:Q}
 Q_i=\sum_{k=0}^{n-2}(k+1)M_{i,0}^{n-2-k}M_{i,1}^k,
 \qquad C_i=\frac{nE_iH^{n-1}}{H^n}.
\end{equation}
The order of the two normal directions is part of these data.

\begin{proposition}\label{intrinsic:intersections}
Let $P$ be the anticanonical polytope of $X$ and $F_i$ its facet
corresponding to $E_i$. On $F_i$, let $t$ be the affine coordinate
which is zero and one on the faces corresponding to $S_{i,0}$ and
$S_{i,1}$, respectively. Then
\begin{align}
 J_i:=\frac{1}{\operatorname{vol}(P)}\int_{F_i}t\,d\operatorname{vol}
 &=\frac{Q_i}{H^n},\label{intrinsic:moment}\\
 C_i&=\frac{n}{H^n}\sum_{k=0}^{n-2}
              M_{i,0}^{n-2-k}M_{i,1}^k.\label{intrinsic:exceptional-degree}
\end{align}
Interchanging the sections gives a positive contribution $J_i^{\rm opp}$
with $J_i+J_i^{\rm opp}=C_i$. If $M_{i,1}-M_{i,0}$ is nef, then
$J_i\ge C_i/2$; if $M_{i,0}-M_{i,1}$ is nef, the inequality reverses.
Numerically equal endpoint classes give $J_i=C_i/2$.
\end{proposition}

\begin{proof}
Let $a_i$ and $c_i$ be the affine support coordinates of the two
original rays defining the $i$-th blowup. The exceptional equation is
$a_i+c_i=1$, with $a_i,c_i\ge0$, so $t=a_i$ has the stated endpoints.
The original two rays extend to an integral basis. On the exceptional
facet, an integral tangent vector can therefore increase $a_i$ by one
and decrease $c_i$ by one. Thus its lattice measure disintegrates as
$dt$ times the lattice measure on the centre, without an index factor.

The slice at $t$ is, up to translation, the moment polytope of
\[
 M_{i,t}=(1-t)M_{i,0}+tM_{i,1}.
\]
One way to see this is to restrict the supporting inequalities of
$H|_{E_i}$ along the two sections: their constants interpolate
affinely, giving precisely the supporting inequalities of $M_{i,t}$.
Both endpoints are ample, so the normal fan remains the fan of $Z_i$
along the entire segment. The slice volume is consequently
$M_{i,t}^{n-2}/(n-2)!$. Since $\operatorname{vol}(P)=H^n/n!$,
\begin{equation}\label{intrinsic:beta-integral}
 J_i=\frac{n(n-1)}{H^n}\int_0^1tM_{i,t}^{n-2}\,dt.
\end{equation}
For $d=n-2$, the coefficient of $M_{i,0}^{d-k}M_{i,1}^k$ in the
integral is
\[
 \binom dk\int_0^1t^{k+1}(1-t)^{d-k}\,dt
 =\frac{k+1}{(d+1)(d+2)}.
\]
This proves \eqref{intrinsic:moment}. Integration without $t$ gives
\eqref{intrinsic:exceptional-degree}, also the projective-bundle
formula for $(H|_{E_i})^{n-1}$. Replacing $t$ by $1-t$ replaces the
coefficient $k+1$ by $n-1-k$, so the two contributions sum to $C_i$.
Their positivity follows from ampleness.

Finally, write $I_k=M_{i,0}^{d-k}M_{i,1}^k$. If
$M_{i,1}-M_{i,0}$ is nef, then
\[
 I_{k+1}-I_k
 =M_{i,0}^{d-k-1}M_{i,1}^k(M_{i,1}-M_{i,0})\ge0.
\]
Pairing $k$ with $d-k$ in
$\sum_k(k-d/2)I_k$ proves $J_i-C_i/2\ge0$.
The reverse nef condition and numerical equality are treated identically.
\end{proof}

\begin{theorem}\label{intrinsic:criterion}
For the blowup and foliation in \eqref{intrinsic:classes},
\begin{equation}\label{intrinsic:defect}
 (b+1)A=b+\sum_iJ_i,\qquad
 A-1=\frac{\sum_iQ_i-H^n}{(b+1)H^n}.
\end{equation}
Consequently $\sum_iQ_i\ge H^n$ excludes stability of $T_X$,
and strict inequality makes $\mathcal F$ a strict destabilizer.
Equivalently, the total contribution is
\begin{align}
 \sum_iJ_i
 &=n\frac{f^*(-K_{\PP^b})H^{n-1}}{H^n}-b\notag\\
 &=(n-b)-n\frac{(-K_{X/\PP^b})H^{n-1}}{H^n}.
 \label{intrinsic:relative-defect}
\end{align}
Here $-K_{X/\PP^b}=H-f^*(-K_{\PP^b})$ denotes the relative
anticanonical divisor class.
\end{theorem}

\begin{proof}
In central affine coordinates the anticanonical polytope has inequalities
\[
 a_j\ge0,\qquad \sum_{j=1}^ba_j\le b+1,\qquad
 a_i+c_i\ge1\quad(1\le i\le m),
\]
together with the inequalities from the rays of $B$. Apply divergence
to $\sum_ja_j\partial/\partial a_j$, whose divergence is $b$.
The coordinate facets and the facets from $B$ have zero flux. The
opposite central facet contributes $(b+1)$ times its lattice volume,
and $F_i$ contributes $-\int_{F_i}a_i$. As in
Lemma~\ref{mixed:flux}, the conversion between Euclidean and lattice
facet measures cancels the normal-length factor. Division by
$\operatorname{vol}(P)$ gives the first identity in
\eqref{intrinsic:defect}. Proposition~\ref{intrinsic:intersections}
gives the second, and Proposition~\ref{intrinsic:foliation} gives the
slope conclusions. Substituting $f^*(-K_{\PP^b})=(b+1)L$ proves
\eqref{intrinsic:relative-defect}.
\end{proof}

Thus an exceptional divisor contributes its normalized degree $C_i$
multiplied by the mean of $t$ in its slice-volume distribution.
The two section polarizations determine how this degree is divided.
For example, the geometric hypotheses $M_{i,1}-M_{i,0}$ nef for every
$i$ and $\sum_iC_i\ge2$ imply failure of stability; a strict weighted
inequality implies instability. We do not impose these nef hypotheses
in the tree theorems below. The identity, which does not require them,
retains the effect of all other blowups through $Z_i$ and $M_{i,0},M_{i,1}$.
It also distinguishes an individual foliation test from a test of every
subsheaf: $\sum_iQ_i<H^n$ controls $\mathcal F$ alone.

\subsection{Geometric slice profiles}

There is a second useful description when $m=b$ and
$B=\prod_{i=1}^bB_i$, with $D_i$ pulled back from a smooth invariant
prime divisor on $B_i$. Put $d_i=\dim B_i$ and $h_i=-K_{B_i}$.
Assume that $h_i$ and $h_i-D_i$ are ample, as well as $H=-K_X$.
Define the positive continuous functions
\begin{equation}\label{intrinsic:profiles}
 p_i(a)=
 \begin{cases}
 (h_i-(1-a)D_i)^{d_i}/d_i!,&0\le a\le1,\\
 h_i^{d_i}/d_i!,&a\ge1.
 \end{cases}
\end{equation}
These functions express the volume of a divisor segment on the base.

\begin{proposition}\label{intrinsic:profile-test}
In this product-base setting, let
\[
 Z(t)=\int_{\substack{a_i\ge0\\\sum_i a_i\le t}}
                 \prod_i p_i(a_i)\,da_1\cdots da_b.
\]
Then
\begin{equation}\label{intrinsic:profile-slope}
 A=\frac{Z'(b+1)}{Z(b+1)},\qquad
 \frac{p_i'(a)}{p_i(a)}
 =\frac{d_iD_i(h_i-(1-a)D_i)^{d_i-1}}
        {(h_i-(1-a)D_i)^{d_i}}\quad(0<a<1).
\end{equation}
If $q_i$ are positive continuous log-concave functions, constant beyond
one, and every $p_i/q_i$ is nonincreasing, then $A$ is at most the
corresponding ratio $Z_q'(b+1)/Z_q(b+1)$. Nondecreasing ratios give
the reverse inequality.
\end{proposition}

\begin{proof}
At fixed $a_i\in[0,1]$, the exceptional constraint moves the facet of
$D_i$ inward by $1-a_i$. The remaining polytope is therefore the
moment polytope of $h_i-(1-a_i)D_i$. The ample endpoint assumption
keeps the whole segment in the same ample cone and gives exactly
\eqref{intrinsic:profiles}. For $a_i\ge1$ the new constraint is
redundant. Independence of the base factors gives the product
integrand. Hence $Z(b+1)$ is the full polytope volume, and $Z'(b+1)$
is the lattice volume of the opposite central facet. Their ratio is
$A$. Differentiating the intersection polynomial gives the second
formula in \eqref{intrinsic:profile-slope}. The comparison is
Lemma~\ref{treemix:convolution}, proved below from Efron's
conditional-sum theorem; it requires log-concavity of the reference
profiles only.
\end{proof}

For a vertex $v$ of a blowup tree, deleting $v$ separates the remaining
construction into factors $B_i$. The divisor at the former parent
incidence becomes an unassigned divisor $D_i$ on its component.
Its class is the pullback of a hyperplane at that component's root.
The classes $h_i-D_i$ are ample: after deleting the parent edge,
every primitive relation of root type has degree at least two.
Subtracting $D_i$ lowers each such degree by one, while edge and
non-root primitive relations are unchanged.
Thus the incoming functions in the tree calculations are precisely the
intersection polynomials \eqref{intrinsic:profiles}.

The proof of Theorem~\ref{treemix:subcubic} bounds their logarithmic
derivatives above by $3/(2+a)$ when every nonleaf dimension equals its
valence and the maximum degree is three. The degree-four obstruction
instead supplies four profiles bounded below in ratio order by
\eqref{treemix:g-four}. Proposition~\ref{intrinsic:profile-test}
allows either comparison to be tested on other polarized base pairs,
without assigning a construction graph to those factors. Outside the
tree criterion, an upper bound concerns only the radial foliation;
a lower bound exceeding one still proves instability of $T_X$.

Neither this description nor \eqref{intrinsic:relative-defect} asserts
that a large Picard number forces the required contraction data or
intersection inequality. That is a separate structural problem. Also,
the edges of the bundle constructions in Section~\ref{sec:graphs}
represent twists, not the exceptional divisors used here; their
subdivision theorems do not assert monotonicity of these blowup weights.
